%--------------------------------------------------------------------------------
%
%
%--------------------------------------------------------------------------------
\pagebreak
\section*{RESET - Reset T64 modules}
\addcontentsline{toc}{section}{RESET - Reset T64 modules}

%--------------------------------------------------------------------------------
\begin{iPageDescEntry}{Syntax}
\texttt{RESET <modNum> | ALL } \\
\end{iPageDescEntry}

%--------------------------------------------------------------------------------
\begin{iPageDescEntry}{Description}

The \texttt{RESET} command resets a module or the entire system. The module will execute its defined reset tasks. The result is module specific. A memory module will clear its memory range, an I/O module will reset its internal state, and a processor will clear its internal registers and start at the architected instruction address.

Like all commands dealing with module management, this command should be taken with care. A common use could be to have this command as the last command in a configuration file to start the Twin-64 system.

\end{iPageDescEntry}

%--------------------------------------------------------------------------------
\begin{iPageDescOpEntry}{Example}
\begin{lstlisting}[style=iPageOpStyle]

RESET 3               # reset module number three.

\end{lstlisting}
\end{iPageDescOpEntry}

%--------------------------------------------------------------------------------
\begin{iPageDescEntry}{Notes}
None.
\end{iPageDescEntry}
